
Prior work on LLM code generation has explored various feedback integration strategies, though systematic comparison of feedback routing policies remains limited.

\textbf{Execution-based feedback systems} such as those described in previous studies have demonstrated iteration reduction through test-driven refinement. These systems typically run execution tests after each generation attempt and provide failure information to guide subsequent iterations.

\textbf{Static analysis integration} has been explored primarily in post-generation validation contexts. Recent work has shown that static analyzers can reduce certain error classes when applied to LLM outputs, though quantitative detection rates in iterative feedback loops have not been systematically measured.

\textbf{Multi-source feedback aggregation} approaches concatenate multiple feedback types simultaneously. No prior work has directly compared sequential routing policies against aggregation baselines with controlled ablation methodology.

The current study differs by isolating the routing mechanism itself: comparing cascade (sequential single-source) against aggregation (simultaneous multi-source) on tasks requiring both type safety and logical correctness validation.
